\section{从“下载互联网”到“下载人类/世界”}

访谈中最有传播力的比喻之一，是“反向 OpenAI”。正向路径是：从互联网下载数据，训练 transformer / GPT，得到语言智能，再推向市场。这个路径之所以成立，是因为互联网已经承载了大量人类写下的知识。世界模型没有同样的 shortcut：真实世界的数据、企业场景、物理过程、机器人任务、传感器反馈，不能简单从网页下载。

\begin{figure}[H]
\centering
\includegraphics[width=0.96\textwidth]{download-internet-to-world.png}
\caption{从“下载互联网”到“下载人类/世界”：世界模型需要真实环境和伙伴闭环。}
\end{figure}

\begin{knowledgebox}{读图：反向路径的关键是合作闭环}
左侧是互联网到语言模型的正向路径；右侧是世界、伙伴、任务和数据共同反哺模型的反向路径。图中下方的闭环条件说明：世界模型需要伙伴数据、初始模型、真实任务、评测反馈和安全控制。没有世界参与，世界模型会停留在论文概念或生成 demo。
\end{knowledgebox}

\subsection{World model needs the world}

谢赛宁说 world model needs the world。这句话是 AMI Labs 组织策略的核心。它意味着世界模型不是一个单公司靠封闭数据就能完成的任务，而需要来自农场、医院、工厂、机器人公司、传感器网络、仿真环境和真实业务场景的合作。

在这个框架里，模型先提供初始能力，进入真实场景创造价值；真实场景产生反馈、错误和新数据；这些数据再反哺模型。这个闭环类似数据引擎，但对象不再只是网页文本，而是物理世界和行动系统。

\begin{importantbox}{反向 OpenAI 的本质}
正向 OpenAI 下载互联网；反向 OpenAI 要和世界共建数据。前者依赖已有文本语料，后者依赖真实任务、伙伴网络、评测体系和持续反馈。
\end{importantbox}

\subsection{为什么这不是单纯买数据}

前面说 world model needs the world，容易让人以为解决方案只是采购更多物理数据。本节要把这个误解拆开：世界模型真正需要的是数据生成机制，而不是一次性数据库存。

如果世界模型需要的是 action 后果、物理状态、失败恢复、传感器流和领域过程，那么“买一批数据”远远不够。数据要有任务、目标、反馈、评测和持续更新。更重要的是，很多数据只有在模型进入场景后才会出现：模型犯错、人工纠正、系统恢复、用户改变流程，这些都不是静态语料库能提前准备好的。

\begin{knowledgebox}{从 data factory 到 environment factory}
本期与 EP134 的数据综述可以连起来看：世界模型需要的不是单次采购的数据，而是能不断产生任务、失败、反馈和评测的环境。未来的数据工厂更像 environment factory。
\end{knowledgebox}

\subsection{本章小结}

“下载互联网”解决了语言智能的第一阶段；“下载人类/世界”意味着从静态文本转向动态环境。世界模型的竞争，不只是模型参数竞争，也是合作网络、评测体系和真实反馈闭环的竞争。

